\documentclass{cmi}
\usepackage{upgreek}
\usepackage[version=4]{mhchem}
\usepackage{subcaption}
\captionsetup[subfigure]{labelformat=simple, labelsep=space}
\renewcommand{\thesubfigure}{\asbuk{subfigure})}

% === БИБЛИОГРАФИЯ BIBLATEX ===
\usepackage[backend=biber, style=gost-numeric, sorting=none]{biblatex}
\addbibresource{references.bib}
% =============================

% Данные поля заполняются редакцией журнала
%\makeatletter
%\newcommand{\@volume}{}
%\newcommand{\@issue}{}
%\newcommand{\@num}{}
%\newcommand{\@pages}{}

\begin{document}
\classify{УДК 519.853.4}

\renewcommand*{\thefootnote}{\fnsymbol{footnote}} %to make a number in end of title delete this line

\title{ПАРАЛЛЕЛЬНЫЙ АЛГОРИТМ ГЛОБАЛЬНОЙ ОПТИМИЗАЦИИ И ЕГО ИСПОЛЬЗОВАНИЕ \\ ДЛЯ РЕШЕНИЯ ОБРАТНЫХ ЗАДАЧ ХИМИЧЕСКОЙ КИНЕТИКИ
\footnote{Статья рекомендована к публикации программным комитетом всероссийской научной конференции <<Параллельные вычислительные технологии (ПаВТ) 2026>>.}%
}

\authors{К.А.~Баркалов$^{1,2}$, Е.Н.~Варсеева$^{1}$, И.Г.~Лебедев$^{1}$, А.Л.~Хашпер$^{3}$, А.С.~Урлуков$^{4,5}$, И.М.~Губайдуллин$^{6}$, И.Ф.~Хисаметдинов$^{7}$}
 
\organizations{$^{1}$Нижегородский государственный университет им. Н.И. Лобачевского \\ (603022 Н.~Новгород, просп.~Гагарина, д.~23)\\
$^{2}$Научно-образовательный математический центр <<Математика технологий будущего>> \\ (603022 Н.~Новгород, просп.~Гагарина, д.~23, корп.~2) \\
$^{3}$ООО <<СИБИНТЕК-СОФТ>> \\ (141401, Московская обл., г. Химки, Коммунальный проезд, стр. 30, помещ. 191)\\
$^{4}$Институт катализа СО РАН \\ (630090, Новосибирск, пр. Академика Лаврентьева, д. 5)\\
$^{5}$Новосибирский государственный университет \\ (630090, Новосибирск, ул. Пирогова, д. 1) \\
$^{6}$Институт нефтехимии и катализа УФИЦ РАН \\ (450075, Уфа, пр. Октября, д. 141)\\
$^{7}$Уфимский государственный нефтяной технический университет \\ (450064, Уфа, ул. Космонавтов, д. 1)\\
}

\newcommand{\emailsinsert}{\href{mailto:konstantin.barkalov@itmm.unn.ru}{konstantin.barkalov@itmm.unn.ru}, \href{mailto:varseeva@unn.ru}{varseeva@unn.ru}, \href{mailto:ilya.lebedev@itmm.unn.ru}{ilya.lebedev@itmm.unn.ru}, \href{mailto: anna.khashper@gmail.com}{ anna.khashper@gmail.com}, \href{mailto: aurlukov@mail.ru}{ aurlukov@mail.ru}, \href{mailto:  irekmars@mail.ru}{ irekmars@mail.ru}, \href{mailto:iskanderkhisametdinov1@gmail.com}{iskanderkhisametdinov1@gmail.com} }

% Дата поступления статьи в редакцию
\newcommand{\rdate}{20.03.2026}
\emails{\emailsinsert}
\redactiondate{\rdate}

\renewcommand*{\thefootnote}{\arabic{footnote}}
\setcounter{footnote}{0}

% оформление аннотации
\begin{abstract}%
В статье рассматривается использование параллельных вычислений для выбора параметров математической модели процесса низкотемпературной паровой конверсии углеводородов, входящих в состав попутного нефтяного газа. Для указанного химического процесса необходимо сформировать его кинетическую модель, то есть определить соответствующие кинетические параметры реакций. Для этого решается обратная задача, в которой ищутся значения кинетических параметров на основе экспериментальных данных. С формальной точки зрения обратная задача химической кинетики является задачей глобальной оптимизации. Для ее решения был использован параллельный информационно-статистический алгоритм глобального поиска в сочетании с локальным уточнением лучшего решения.
Алгоритм является детерминированным и основан на предположении липшицевости целевой функции, что является типичным для многих других подходов к построению методов глобальной оптимизации. При этом решение многомерных задач сводится к решению эквивалентных им одномерных задач. Соответствующая редукция основана на использовании кривой Пеано, отображающей единичный отрезок вещественной оси на гиперкуб. Распараллеливание алгоритма организовано по схеме «мастер-рабочие» с использованием общей памяти. 
Оптимальные параметры модели, найденные с помощью параллельного алгоритма оптимизации, позволили адекватно смоделировать процесс низкотемпературной паровой конверсии легких углеводородов с использованием катализатора.

\keywords{глобальная оптимизация, многоэкстремальные функции, параллельные вычисления, химическая кинетика, обратные задачи}
\end{abstract}
\newpage
%Образец цитирования статьи
\forcitation{Баркалов~К.А., Варсеева~Е.Н., Лебедев~И.Г., Хашпер~А.Л., Урлуков~А.С., Губайдуллин~И.М., Хисаметдинов~И.Ф. Параллельный алгоритм глобальной оптимизации и его использование для решения обратных задач химической кинетики}

%Верхние колонтитулы
\markboth{К.А.~Баркалов, Е.Н.~Варсеева и др.}{Параллельный алгоритм глобальной оптимизации\ldots}

% Заголовки разделов формируются при помощи команд  \section{}, \subsection{}, \subsubsection{}
\section*{Введение}
\label{sec-intro}

%Первый вариант введения
%В настоящее время сжигание попутного нефтяного газа (ПНГ) является актуальной проблемой российской нефтедобывающей промышленности. ПНГ представляет собой смесь предельных углеводородов, добываемых на нефтяных месторождениях. Одним из наиболее перспективных способов эффективной утилизации ПНГ является использование его непосредственно на месторождении для выработки электроэнергии и тепла на малых электростанциях для обеспечения собственных нужд месторождения~\cite{UskovPSSKS2019}. Использование исходного ПНГ в качестве топлива неприемлемо, поскольку он имеет высокое содержание $C_{2+}$-углеводородов, что может вызывать поломку двигателей внутреннего сгорания. Для решения этой проблемы предлагается использовать низкотемпературную паровую конверсию (НТПК) $C_{2+}$-углеводородов, входящих в состав ПНГ, с преимущественным образованием метан-водородных смесей~\cite{ZyryanovaSSBKS2014}, которые уже можно использовать в качестве топлива.

%Для проведения НТПК $C_{2+}$-углеводородов были широко исследованы катализаторы на основе никеля~\cite{EnikeevaPUSSEG21,PotemkinUGKSB2020,ShigarovUP2022,UskovPSSKS2019}. Однако основной проблемой никелевых катализаторов остаётся их невысокая стабильность, связанная с зауглероживанием или спеканием никелевых частиц во время эксплуатации катализатора.

%Сотрудники Института катализа СО РАН в своих работах показали, что родиевые катализаторы, нанесённые на смешанные оксиды церия-циркония, проявляют высокую активность и стабильность в процессе НТПК пропана~\cite{UrlukovUSP22,UrlukovUSBRRS23,GarkulZFBBP2024,UrlukovUSSP2022,ZadesenetsGFSPP2023}, а катализатор 1~вес.\,\%~Rh/Ce$_{0.4}$Zr$_{0.5}$Y$_{0.05}$La$_{0.05}$O$_2$ является перспективным вариантом среди других родий-содержащих аналогов~\cite{UrlukovUSBRRS23}. Однако для более детального понимания кинетики процесса НТПК углеводородов на Rh катализаторах и оптимизации условий его проведения нужно построить и проанализировать математическую модель указанного процесса (так называемую кинетическую модель).

%Обновленный вариант введения

Сжигание попутного нефтяного газа (ПНГ) на текущий момент остается актуальной проблемой российской нефтедобывающей промышленности. 
ПНГ представляет собой смесь предельных углеводородов, растворенных в нефти и выделяющихся при ее добыче. 
Одним из наиболее перспективных способов эффективной утилизации попутного газа является его использование непосредственно на месте добычи, для выработки электроэнергии и тепла на малых электростанциях, которые обеспечивают собственные нужды месторождения~\cite{UskovPSSKS2019}. Однако исходный ПНГ не может быть использован в качестве топлива из-за высокого содержания в нем \ce{C2+}-углеводородов, которые вызывают поломки двигателей внутреннего сгорания. Для решения проблемы утилизации ПНГ предлагается применять низкотемпературную паровую конверсию (НТПК) \ce{C2+}-углеводородов, входящих в состав ПНГ, с преимущественным образованием метан-водородных смесей~\cite{ZyryanovaSSBKS2014}. Полученные таким образом метан-водородные смеси могут быть использованы в качестве топлива.

Для проведения НТПК \ce{C2+}-углеводородов были широко исследованы катализаторы на основе никеля~\cite{EnikeevaPUSSEG21,PotemkinUGKSB2020,ShigarovUP2022,UskovPSSKS2019}. Однако основной проблемой никелевых катализаторов является их невысокая стабильность, связанная с зауглероживанием или спеканием никелевых частиц во время эксплуатации катализатора.

Сотрудники Института катализа СО РАН в своих работах показали, что родиевые катализаторы, нанесенные на смешанные оксиды церия-циркония, проявляют высокую активность и стабильность в процессе НТПК пропана~\cite{UrlukovUSP22,UrlukovUSBRRS23,GarkulZFBBP2024,UrlukovUSSP2022,ZadesenetsGFSPP2023}, а катализатор 1~вес.\,\%~Rh/Ce$_{0.4}$Zr$_{0.5}$Y$_{0.05}$La$_{0.05}$O$_2$ является оптимальным вариантом среди родий-содержащих катализаторов~\cite{UrlukovUSBRRS23}.
Однако для более детального понимания кинетики процесса НТПК углеводородов на Rh катализаторах и оптимизации условий его проведения нужно построить и проанализировать математическую модель указанного процесса (так называемую кинетическую модель).

%Кинетическая модель является основой математического моделирования химических реакций и исходным уровнем при математическом моделировании химических реакторов. 
%При разработке кинетической модели многостадийной химической реакции требуется решить обратную задачу химической кинетики, при решении которой на основе экспериментальных данных о динамике протекания процесса определеляются количественные характеристики кинетической модели --- константы скоростей одностадийных реакций. С математической точки зрения обратная задача химической кинетики представляет собой задачу глобальной оптимизации, поэтому важно выбрать подходящий для ее решения метод.

При построении кинетической модели многостадийной химической реакции решается обратная задача химической кинетики: по экспериментальным данным о динамике процесса восстанавливаются численные параметры модели --- константы скоростей одностадийных реакций. Формально все сводится к поиску глобального минимума многоэкстремальной целевой функции, поэтому очень важно использовать адекватный метод оптимизации.

%Смысловая связка между фрагментами
%Метод решения оптимизационных задач подобного типа должен учитывать следующую специфику. Во-первых, целевая функция задачи представляет собой <<черный ящик>>, поскольку она задается не аналитической формулой, а в виде алгоритма, вычисляющего ее значения. Во-вторых, даже однократное вычисление значения целевой функции является трудоемкой операцией, так как требует выполнения численного моделирования. Указанная специфика исключает возможность применения для решения обратных задач мультистартовых методов, основанных на многократном использовании градиентного спуска. С одной стороны, вычисление градиента в данном случае оказывается чрезмерно дорогостоящим, а с другой --- мультистартовые подходы слабо применимы к существенно многоэкстремальным задачам, каковыми и являются обратные задачи химической кинетики.

Метод решения таких оптимизационных задач должен учитывать их специфику. Во-первых, целевая функция фактически представляет собой <<черный ящик>>: она задается не явной формулой, а алгоритмом, вычисляющим ее значения. Во-вторых, даже одно вычисление функции требует ресурсов, поскольку подразумевает проведение численного моделирования. Из-за этого мультистартовые методы, основанные на многократном применении градиентного спуска, неприменимы: вычисление градиента слишком дорого, и даже многократный запуск градиентных процедур может не обеспечить сходимость к решению многоэкстремальной задачи.

Перечисленные особенности обратных задач учитываются в разработанном в Нижегородском университете параллельном алгоритме глобальной оптимизации. Теоретической основой данного алгоритма служит информационно-статистический подход, подробно изложенный в \cite{Strongin2000}. Суть подхода заключается в редукции исходной многомерной задачи к эквивалентной одномерной задаче, которую можно решить с помощью эффективного алгоритма оптимизации функций одной переменной. Вопросы распараллеливания предложенного алгоритма на различных архитектурах рассматриваются в \cite{Barkalov2016}.

%Структура статьи
В настоящей статье представлены результаты применения параллельного алгоритма глобального поиска к решению обратной задачи химической кинетики. Статья имеет следующую структуру. 
В разделе \ref{sec-problem} приведено описание математической модели изучаемой химической реакции, сформулирована обратная задача.
В разделе \ref{sec-PGSA} излагается стандартная постановка задачи глобальной оптимизации, приведена вычислительная схема параллельного алгоритма ее решения, сформулированы его теоретические свойства, кратко рассмотрены вопросы, связанные с техническими аспектами реализации параллельного алгоритма.
Содержательная интерпретация результатов численных расчетов с использованием разработанного параллельного алгоритма представлена в разделе \ref{sec-result}.
В заключении содержатся выводы и планы дальнейших исследований.

\section{Постановка задачи}
\label{sec-problem}

Для описания низкотемпературной паровой конверсии пропана достаточно учесть следующие реакции:

\begin{equation}\label{eq1}
\begin{array}{l}
\ce{C3H8 + 6H2O -> 3CO2 + 10H2}, \\
\ce{C3H8 + H2 -> C2H6 + CH4}, \\
\ce{C2H6 + H2 -> 2CH4}, \\
\ce{CO2 + 4H2 <=> CH4 + 2H2O}. \\
\end{array}
\end{equation}

Было проведено математическое моделирование НТПК этана и пропана на 1~вес.\,\%~Rh/Ce$_{0.4}$Zr$_{0.5}$Y$_{0.05}$La$_{0.05}$O$_{2}$ катализаторе с использованием реакций~(1) и механизма Ленгмюра--Хиншельвуда для задания скоростей некоторых реакций~\cite{ArcotumapathyAAVAA2015,BatebiAM2021}:
\begin{equation}
\begin{cases}
\displaystyle w_{1} = k_{1} \cdot \frac{P_{\ce{C3H8}}\cdot P_{\ce{H2O}}}{(P_{\ce{H2O}}+\alpha\cdot\sqrt{P_{\ce{H2}}})^2}, \\
\displaystyle w_{2} = k_{2} \cdot P_{\ce{C3H8}} \cdot P_{\ce{H2}}, \\
\displaystyle w_{3} = k_{4}\cdot \frac{\sqrt{P_{\ce{H2}}}\cdot P_{\ce{C2H6}}}{(P_{\ce{H2O}}+\alpha\cdot\sqrt{P_{\ce{H2}}})^2},\\
\displaystyle w_{4} = k_{4}\cdot P_{\ce{CO2}}\cdot P_{\ce{H2}}-k_{-4}\cdot P_{\ce{H2O}}\cdot P_{\ce{CH4}},\\
\end{cases}
\end{equation}
где \textit{P}$_{X}$ --- парциальное давление вещества $X$, атм; \textit{w}$_{i}$ --- скорость $i$-ой реакции, атм/(л$\cdot$с); $k_{i}$ --- константа скорости $i$-ой реакции; $t$ --- время, с; $\alpha$ --- отношение констант адсорбции \ce{H2} и \ce{H2O} на выбранном катализаторе.

Система дифференциальных уравнений для описания кинетики НТПК пропана записывается в виде
\begin{equation}
\begin{cases}
\displaystyle u\cdot \frac{dP_{\ce{C3H8}}}{dl} = -w_{1}-w_{2},\\
\displaystyle u\cdot \frac{dP_{\ce{H2O}}}{dl} = -6w_{1}+2w_{4},\\
\displaystyle u\cdot \frac{dP_{\ce{CO2}}}{dl} = 3w_{1}-w_{4},\\
\displaystyle u\cdot \frac{dP_{\ce{H2}}}{dl} = 10w_{1}-w_{2}-w_{3}-4w_{4},\\
\displaystyle u\cdot \frac{dP_{\ce{CH4}}}{dl} = w_{2}+2w_{3}+w_{4},\\
\displaystyle u\cdot \frac{dP_{\ce{C2H6}}}{dl} = w_{2}-w_{3},\\
\end{cases}
\end{equation}
где $u$ --- линейная скорость потока реакционной смеси, м/с; $l$ --- высота слоя катализатора, м.

Прямая задача химической кинетики предполагает расчет состава многокомпонентной реагирующей смеси на основе кинетической модели с известными параметрами~\cite{MazitovOAG16}.
С математической точки зрения прямая задача представляет собой задачу Коши для системы из $n$ обыкновенных дифференциальных уравнений:
\begin{equation}
x'=f(x, t), x(t_{0}) = x_{0}, t\in[t_{0},t_{cont}],
\end{equation}
где $x(t)$ --- вещественная $n$-мерная функция, представляющая собой парциальные давления веществ; $x_{0}$ --- вектор начальных условий; $t_{0}$ --- начальный момент времени; $t_{cont}$ --- время контакта с катализатором; $t \in [t_{0},t_{cont}]$ --- время (независимая переменная).
Начальное и конечное значения $t_{0}$ и $t_{cont}$, а также исходные парциальные давления веществ $x_{0}$ считаются заданными.
Решением прямой задачи является набор $(t_i, x_i)$, где $t_i$ --- момент времени, а $x_i$ --- вектор парциальных давлений всех веществ системы в $i$-й момент времени~\cite{MazitovOAG16}.

Для решения прямой задачи при различных температурах реакции был реализован модифицированный метод Рунге-Кутты 4 порядка точности в соответствии со следующими формулами:

\begin{equation}
\begin{array}{l}
F_{1}=f(x_{i})\cdot h,\\
F_{2}=f(x_{i}+F_{1}/3)\cdot h,\\
F_{3}=f(x_{i}+(F_{1}+F_{2})/6)\cdot h,\\
F_{4}=f(x_{i}+(F_{1}+3 \cdot F_{3})/8)\cdot h,\\
x_{s+1}=x_{i} + (F_{1}-3 \cdot F_{3} + 4 \cdot F_{4})/ 2,\\
F_{5}=f(x_{i+1})\cdot h,\\
x^*_{i+1}=x_{i} + (F_{1}+4 \cdot F_{4} + F_{5})/ 6.\\
\end{array}
\end{equation}

Построение математических моделей химических реакций предполагает наличие в них неизвестных кинетических параметров (предэкспоненциальные множители констант скоростей, энергии активации, наблюдаемые порядки по компонентам), которые можно найти, решая задачу минимизации отклонения между расчетными и экспериментальными данными. Таким образом, возникает задача идентификации математической модели (обратная задача химической кинетики)~\cite{GubaidullinRT11,MorlanesLYMC2022,WooPBP2023}.

Решение обратной задачи химической кинетики сводится к поиску минимума целевой функции~\cite{MazitovOAG16}. В данной работе в качестве критерия минимизации была использована целевая функция вида
\begin{equation}\label{eq6}
E=\sum_{i=1}^n \left|\overline{P_{i}}-\widetilde{P_{i}}\right|\cdot v_{i},
\end{equation}
где $\overline{P_{i}}$, $\widetilde{P_{i}}$ --- экспериментальные и расчетные значения парциальных давлений веществ, атм; $v_{i}$ --- весовые коэффициенты; $n$ --- количество измеряемых веществ.
Требуется найти набор констант скоростей $k$, при котором значение функции (\ref{eq6}) минимально.
%Содержательная постановка задачи. Математическая модель в виде системы ОДУ. Метод решения ОДУ, его свойства. Задача идентификации параметров реакции.

\section{Параллельный алгоритм для решения задач глобальной оптимизации} \label{sec-PGSA}
\subsection{Задача глобальной оптимизации и редукция размерности}

Рассмотрим задачу глобальной оптимизации в общем виде:
\begin{gather}
 \varphi(y^\ast)=\min{\left\{\varphi(y):y\in D\right\}}, \label{problem}\\
 D=\left\{y\in \mathbb{R}^N: a_i\leqslant y_i \leqslant b_i, \; a_i,b_i\in \mathbb{R}, \;  1\leqslant i \leqslant N\right\} \label{D},
\end{gather}
где целевая функция является <<черным ящиком>> и предполагается удовлетворяющей условию Липшица
\[
\left|\varphi(y_1)-\varphi(y_2)\right|\leqslant L\left\|y_1-y_2\right\|,\; y_1,y_2 \in D,
\]
с константой $L, \; L<\infty,$ неизвестной априори.

Предположение о выполнении условия Липшица для целевой функции используется многими исследователями при разработке детерминированных алгоритмов глобальной оптимизации.
Первые методы такого класса были предложены еще в начале 1970-х годов \cite{Piyavskii1972,Shubert1972}. Активная работа в этой области продолжается и сейчас \cite{Evtushenko2013,Zilinskas2010,Pinter1996,Jones2009}.

%В настоящее время для решения задач с целевой функцией вида <<черный ящик>> широко применяются биоинспирированные метаэвристические алгоритмы; см., например, \cite{Yang2013,Gendreau2010,Eiben2015}. Такие алгоритмы так или иначе используют идеи случайного поиска. Благодаря простоте реализации и применения они получили широкое распространение. Однако по качеству работы (например, по количеству корректно решенных задач из заданного класса) эти алгоритмы уступают детерминированным аналогам \cite{Kvasov2018,Sergeyev2018}.

Отметим, что в настоящее время распространенным инструментом для решения задач с функциями вида <<черный ящик>> являются также и метаэвристические алгоритмы, основанные на имитации природных процессов и использующие элементы стохастики (см. \cite{Yang2013,Gendreau2010,Eiben2015}). Их практическая привлекательность объясняется простотой использования, однако по эффективности и надежности поиска оптимума они, как правило, уступают детерминированным методам \cite{Kvasov2018,Sergeyev2018}.


%Одним из эффективных детерминированных методов решения задач многоэкстремальной оптимизации является \textit{информационно-статистический алгоритм глобального поиска} \cite{Strongin2000}. Первоначально предложенный для задач без ограничений, он был успешно обобщен на классы задач с невыпуклыми ограничениями и многокритериальной оптимизации. Для различных версий алгоритма также были предложены методы распараллеливания, учитывающие архитектуру современных вычислительных систем \cite{Barkalov2016,globalizerSystem,Strongin2018}.

Эффективный алгоритм решения задач многоэкстремальной оптимизации --- \textit{информационно-статистический алгоритм глобального поиска} --- был разработан в ННГУ им. Н.И. Лобачевского \cite{Strongin2000}. 
Данный алгоритм является детерминированным и в исходном виде был ориентирован на решение задач безусловной оптимизации, однако позднее были разработаны его модификации для решения задач с функциональными ограничениями, а также --- с несколькими критериями. 
Также для алгоритма глобального поиска был предложен ряд способов его распараллеливания \cite{Barkalov2016,globalizerSystem,Strongin2018}.

Многие современные алгоритмы глобальной оптимизации используют идею редукции размерности при решении многомерных задач; см., например, методы на основе диагональных \cite{Sergeyev2006} или симплексных \cite{Zilinskas2008} разбиений. Параллельный алгоритм глобального поиска, описываемый в данном разделе, использует кривую Пеано \cite{Sergeyev2013,Strongin2000}, отображающую единичный отрезок $[0,1]$ на $N$-мерный куб $D$ из (\ref{D}). С помощью такого отображения исходная многомерная задача (\ref{problem}) сводится к одномерной задаче
\[
\varphi(y^\ast)=\varphi(y(x^\ast))=\min{\left\{\varphi(y(x)): x\in[0,1]\right\}},
\]
где функция $\varphi(y(x))$ удовлетворяет равномерному условию Гёльдера
\[
\left|\varphi(y(x_1))-\varphi(y(x_2))\right|\leqslant H\left|x_1-x_2\right|^{1/N}
\]
с константой Гёльдера $H$, связанной с константой Липшица $L$ соотношением 
$ H=2 L \sqrt{N+3}$, а $y(x)$ --- кривая Пеано, отображающая $[0,1]$ на $D$.
Поэтому, не ограничивая общности, можно говорить о минимизации одномерной функции
\[
f(x) = \varphi(y(x)), x\in[0,1],
\]
удовлетворяющей условию Гёльдера.
%комментарий про гельдеовскую функцию f(x)

Известно, что точная кривая Пеано $y(x)$ задается как предельный объект и вычислить точное значение $y(x)$ не представляется возможным. Поэтому на практике используются аппроксимации истинной кривой, называемые \textit{развертками}. Численные методы построения разверток изложены в работах \cite{Sergeyev2013,Strongin2000}. Точность развертки определяется параметром ее построения $m$ и имеет порядок $2^{-m}$ по каждой координате.
Примеры разверток для размерности $N=3$ и различных значений $m$ показаны на рис.~\ref{evolvents}.

\begin{figure}[t]
\centering
\begin{subfigure}[b]{0.48\textwidth}
    \centering
    
\includegraphics[width=\linewidth]{fig1a.JPG}
    \caption{Развертка при $m=3$}
    \label{fig:evolvent_m3}
\end{subfigure}
\hfill
\begin{subfigure}[b]{0.48\textwidth}
    \centering
    \includegraphics[width=\linewidth]{fig1b.JPG}
    \caption{Развертка при $m=4$}
    \label{fig:evolvent_m4}
\end{subfigure}
\caption{Развертки в трехмерном пространстве}
\label{evolvents}
\end{figure}

\subsection{Параллельный алгоритм глобального поиска}

Будем называть процесс вычисления значения функции $f(x)=\varphi(y(x))$ \textit{испытанием}, а пару $\{x, z = f(x) = \varphi(y(x))\}$ --- результатом испытания.
При описании параллельного алгоритма будем использовать термин \textit{итерация} для обозначения одновременного (параллельного) выполнения нескольких испытаний (каждое испытание выполняется в отдельном потоке или процессе). Число испытаний на $n$-й итерации обозначим $p$, а общее число испытаний, выполненных за все $n$ итераций, обозначим $s(n)$. Иными словами, предполагается, что при выполнении $n$-й итерации имеется в распоряжении $p$ потоков или процессов.

Вычислительная схема параллельного алгоритма глобального поиска состоит в следующем.
Первые $p$ испытаний выполняются в граничных точках $x^0 = 0$, $x^1 = 1$, а также в произвольных внутренних точках $x^2, ..., x^{p-1}$ интервала $(0,1)$. Предположим, что завершено $n \geqslant 1$ итераций метода, в ходе которых проведены испытания в $s=s(n)$ точках $x^i, 0 \leqslant i \leqslant s$. Тогда точки $x^{s+1},...,x^{s+p}$ для проведения испытаний на следующей $(n+1)$-й итерации определяются согласно следующим правилам. 

\begin{enumerate}
	\item 
	Перенумеровать прообразы всех точек, в которых уже проведены испытания,  
\begin{equation}\label{y_i}
y^0=y(x^0), y^1=y(x^1),...,y^s=y(x^s)
\end{equation}
по возрастанию их координат, то есть
\begin{equation}\label{x_i}
0=x_0<x_1<\dots <x_s=1,
\end{equation}
и сопоставить им вычисленные значения $z_i=\varphi(y(x_i)), 0\leqslant i \leqslant s$.

\item
Вычислить максимальное абсолютное значение разделенных разностей
\begin{equation}\label{mu}
\mu = \max_{1 \leqslant i \leqslant s}\frac{\left|z_i-z_{i-1}\right|}{\Delta_i},
\end{equation}
где $\Delta_i=\left(x_i-x_{i-1}\right)^{1/N}$. Если $\mu = 0$, полагаем $\mu = 1$.

\item
Для каждого интервала $(x_{i-1}, x_i), \; 1\leqslant i \leqslant s,$ вычислить величину
\begin{equation}\label{R}
R(i)=r\mu\Delta_i+\frac{(z_i-z_{i-1})^2}{r\mu\Delta_i}-2(z_i+z_{i-1})
\end{equation}
называемую \textit{характеристикой} интервала; вещественное число $r>1$ является параметром алгоритма.

\item 
Упорядочить характеристики $R(i), 1 \leqslant i \leqslant s$, по убыванию 
\begin{equation}\label{Rdec}
R(t_1)\geqslant R(t_2)\geqslant \dots \geqslant R(t_{s})
\end{equation}
и выбрать $p$ интервалов с номерами $t_j, 1\leqslant j \leqslant p$, которым соответствуют наибольшие характеристики.

\item
Провести новые испытания параллельно, в точках $x^{s+j}\in(x_{t_j-1},x_{t_j}), 1\leqslant j\leqslant p$, вычисляемых по формуле
\[
x^{s+j} = \frac{x_{t_j}+x_{t_j-1}}{2} - \mathrm{sign}(z_{t_j}-z_{t_j-1})\frac{1}{2r}\left[\frac{\left|z_{t_j}-z_{t_j-1}\right|}{\mu}\right]^N.
\]

\end{enumerate}

Алгоритм завершает свою работу, если условие $\Delta_{t_j}<\epsilon$ выполняется хотя бы для одного номера интервала $t_j, 1 \leqslant j \leqslant p$; здесь $\epsilon>0$ --- заданная точность. Алгоритм также останавливается, если проведено максимально допустимое число итераций $S_{max}$, установленное до начала вычислений.

%Теоретические сведения о распараллеливании
Введение описанного выше параллелизма не приводит к возникновению новых предельных точек, отличных от предельных точек последовательного алгоритма, т.е. когда $p = 1$. Условия сходимости для параллельного АГП могут быть сформулированы в виде следующей теоремы, доказательство которой приведено в работе \cite{Strongin2000}.

\begin{theorem*}%
	Пусть $\overline{x}$ есть предельная точка последовательности испытаний $\{x^s\}$, порождаемой параллельным АГП при минимизации функции $f(x), x\in [0,1],$ удовлетворяющей условию Гёльдера с константой $H, 0<H<\infty$; число задействованных вычислительных элементов $p, 1<p<\infty,$ является фиксированным, а в условии остановки используется значение $\epsilon = 0$. Тогда:
	\begin{enumerate}
		\item
		сходимость к точке $\overline{x} \in (0,1) $ является двусторонней;
		\item
		точка $\overline{x}$ является локально оптимальной, если функция $f(x)$ имеет конечное число локальных экстремумов;
		\item
		если наряду с точкой $\overline{x}$ существует другая предельная точка $\widehat{x}$, то $f(\overline{x})=f(\widehat{x})$;
		\item
		при любом $s>1$ выполняется неравенство $f(x_s) \geqslant f(\overline{x})$;
		\item
		если на некотором шаге для величины $r\mu$ из (\ref{R}) выполняется условие $r\mu > 2^{2-1/N}H$;
	\end{enumerate}
	то множество предельных точек последовательности $\{x^s\}$ совпадает с множеством точек глобального минимума $f(x)$.
\end{theorem*}

%Комментарий о распараллеливании
Рассматриваемый вариант распараллеливания является \textit{синхронным}, когда переход к следующей итерации осуществляется только после полного завершения текущей, то есть после окончания последнего испытания текущей итерации. 
Конкретная реализация этой схемы распараллеливания будет зависеть от архитектуры компьютерной системы, а также от требований к программному обеспечению, необходимому для вычисления значений целевой функции. 
В рассматриваемой задаче химической кинетики для проведения одного испытания требовалось решить систему ОДУ (для решения этой задачи использовался модифицированный метод Рунге-Кутты 4-го порядка). Распараллеливание было организовано с использованием общей памяти и схемы мастер-рабочие. Алгоритм глобального поиска был запущен в потоке-мастере, остальные (подчиненные) потоки проводили параллельные испытания в рамках шага 5 алгоритма. При этом время проведения одного испытания составляло порядка 0.1 секунды, таким образом, накладные расходы на синхронизацию между потоками не оказывали существенного влияния на эффективность распараллеливания.

%Комментарий о локальном уточнении
После определения с помощью параллельного АГП точки, являющейся грубой оценкой глобального минимума, из нее запускается локальный метод. Мы использовали метод Хука--Дживса, который относится к классу методов нулевого порядка. Его вычислительные правила основаны на чередовании исследующего поиска (для выбора направления) и поиска в выбранном направлении (поиска по образцу) \cite{Himmelblau75}.

Исследующий поиск можно рассматривать как оператор $y^s=F(w^s)$, отображающий текущую точку $w^s$ в следующую базовую точку $y^s$:
\begin{itemize} 
  \item Вычисляется величина смещения из текущей точки $w^s$ (она является различной для каждого координатного направления и может изменяться в процессе работы метода). 
  \item Шаг поиска является успешным, если значение целевой функции в новой точке не превышает значение целевой функции в предыдущей точке. 
  \item Если шаг является неуспешным, то происходит возврат к исходной точке и выполняется смещение в обратном направлении. 
  \item После исследования всех $N$ координатных направлений данный этап завершается; полученная точка $y^s$ определяется как базовая для поиска по направлению.
\end{itemize}

Поиск по направлению определяется смещением от полученной новой базовой точки $y^s$ вдоль прямой, соединяющей $y^s$ и $y^{s-1}$, т.е.
\[
w^{s+1}=y^s + \alpha (y^s-y^{s-1}),
\]
где величина смещения задается параметром метода $\alpha$.

Рис. \ref{fig:fig1} иллюстрирует работу локального алгоритма. Показаны линии уровня целевой функции; темные точки соответствуют успешным, а светлые --- неуспешным шагам, проведенным в ходе исследовательского поиска.

\begin{figure} 
	\begin{center} 
			\includegraphics[width=0.8\linewidth]{pic/fig1.pdf} 
			\caption{Иллюстрация работы метода Хука-Дживса} %%  подпись к рисунку 
			\label{fig:fig1} 
	\end{center} 
\end{figure}	



\section{Результаты численных экспериментов}
\label{sec-result}

Вычислительные эксперименты проводились на узле суперкомпьютера <<Лобачевский>> в следующей конфигурации: операционная система Linux CentOS 8.2, 2 процессора Intel Xeon Platinum 8558 2.1/4GHz (48 ядер каждый), оперативная память 2048 Gb. 
Все рассмотренные алгоритмы реализованы на С++, использовался компилятор Intel С++ Сompiler 2023 с поддержкой OpenMP.

Выбор оптимальных кинетических параметров осуществлялся следующим образом. Были получены экспериментальные данные для 18 различных температур проведения реакции. Далее, для каждой из этих температур решалась обратная задача и выбирались константы скоростей реакций, минимизирующие функцию (\ref{eq6}) отклонения расчетных концентраций от экспериментальных для каждой температуры. Далее предэкспоненциальный множитель и энергии активации, входящие в уравнения линейно, искались с помощью метода наименьших квадратов.

Таким образом, с помощью параллельного АГП было решено 18 обратных задач, каждая для своей температуры.
Параллельный алгоритм запускался со следующими параметрами: плотность построения развертки $m=10$, параметр надежности $r=1.5$, точность глобального поиска $\epsilon = 10^{-2}$, локального уточнения $\delta = 10^{-4}$; максимальное число итераций глобального поиска $S_{glob} = 64 000$, локального поиска $S_{loc} = 500$. Число потоков варьировалось от 1 до 96.
В табл.~\ref{tab:scalability} и на рис.~\ref{fig:performance_plots} приведены усредненные (по всем 18 задачам) оценки времени работы и ускорения параллельного алгоритма. Результаты показывают, что увеличение числа задействованных потоков (более 32) не приводит к значительному ускорению; этот факт обсуждается в заключении.
Подробный анализ и сопоставление теоретического (по закону Амдала) и практического ускорения параллельного алгоритма глобального поиска с использованием локального уточнения проведен в работе \cite{Lebedev2022}.

\begin{table}[h]
\noindent
\caption{Зависимость времени решения и ускорения от числа потоков}
\label{tab:scalability}
\centering
\begin{tabular}{|c|c|c|}
\hline
\textbf{P} & \textbf{Время решения, с} & \textbf{Ускорение} \\
\hline
1   & 1170.00  & --- \\
2   & 600.27  & 1.95 \\
4   & 308.34  & 3.79 \\
8   & 162.82  & 7.19 \\
16  & 87.95   & 13.30 \\
32  & 54.43   & 21.50 \\
64  & 50.80   & 23.20 \\
96  & 47.90   & 24.42 \\
\hline
\end{tabular}
\end{table}
%Для наглядного представления экспериментальных данных, приведенных в таблице~\ref{tab:scalability}, на рис.~\ref{fig:performance_plots} построена диаграмма зависимости ускорения решения параллельного алгоритма от числа процессов. 

\begin{figure}[h]
\centering
\begin{subfigure}[c]{0.48\textwidth}
    \centering
    \includegraphics[width=\linewidth, height=5cm]{t_N1.png}
    \caption{зависимость времени от числа потоков}
    \label{fig:time_plot}
\end{subfigure}
\hfill
\begin{subfigure}[c]{0.48\textwidth}
    \centering
    \includegraphics[width=\linewidth, height=5cm]{speedup_plot_N1.png}
    \caption{зависимость ускорения от числа потоков}
    \label{fig:speedup_plot}
\end{subfigure}
\caption{Графики времени работы и ускорения параллельного алгоритма}
\label{fig:performance_plots}
\end{figure}
С помощью алгоритма глобального поиска для каждой температуры были получены наборы констант скоростей, минимизирующие функцию отклонения расчетных концентраций от экспериментальных (\ref{eq6}) для каждой температуры. В качестве весовых коэффициентов $v_i$ в формуле (\ref{eq6}) были использованы значения $v_i= 1.0, 1\leqslant i\leqslant 5, v_6=4.4$, что соответствует большему влиянию последнего вещества. В ходе численных экспериментов было установлено, что равномерное взвешивание всех компонентов ($v_{i} = 1$) приводит к систематическому завышению расчётных значений для \ce{C2H6}. Увеличение его весового коэффициента до 4.4 позволило скорректировать эту погрешность и добиться более равномерного распределения ошибки по всем компонентам.

После того, как с помощью параллельного АГП были получены наборы констант скоростей для каждой температуры, осуществлялся поиск предэкспоненциального множителя $k_0$ и энергии активации $E_a$ (Дж/моль). Для этого метод наименьших квадратов применялся к соотношению, полученному логарифмированием уравнения Аррениуса:
\begin{equation}
\ln (k)=\ln (k_0) - \frac{E_a}{RT},
\end{equation}
где $R$ --- универсальная газовая постоянная, Дж/(моль $\cdot$ K); $T$ --- температура, $K$.
Здесь в качестве $k$ подставлялись полученные параллельным алгоритмом глобального поиска значения констант скоростей, соответствующие различным температурам $T$.

Далее методом наименьших квадратов были рассчитаны величины $k_0$ и $E_a$ для всех стадий; полученные значения приведены в табл.~\ref{tab:kinetik}.

\begin{table}[h]
\noindent
\caption{Рассчитанные кинетические параметры реакций (\ref{eq1})}
\label{tab:kinetik}
\centering
\begin{tabular}{|c|c|c|}
\hline
\textbf{№ реакции} & \textbf{$ln(k_0)$} & \textbf{$E_a$, кДж/моль} \\
\hline
1   & 15.09  & 73.28 \\
2   & 4.65   & 7.51 \\
3   & 3.19   & 8.88 \\
4   & 4.28   & 5.87 \\
4'  & 10.6   & 57.31 \\
\hline
\end{tabular}
\end{table}
Полученные значения энергии активации были сопоставлены со значениями, представленными в других работах по исследованию процесса паровой конверсии углеводородов на родиевых катализаторах.
В работе \cite{Schadel2009} был изучен паровой риформинг метана, этана, пропана и бутана на монолитном сотовом катализаторе на основе родия. Эффективные энергии активации паровой конверсии пропана и этана в процессе парового риформинга природного газа составляют 51 и 50 кДж/моль соответственно.
В работе \cite{Karakaya2017} был представлен механизм поверхностной реакции частичного окисления и паровой конверсии пропана на катализаторе Rh/Al$_{2}$O$_{3}$. Процессы разрыва C-C и C-H связей в пропане характеризуются невысокими энергиями активации в пределах 20--70 кДж/моль.

Экспериментальные исследования реакции метанирования CO$_{2}$ на родиевых катализаторах \cite{Solymosi1981} показали, что энергия аткивации этого процесса варьируется от 16.2 до 24.0 ккал/моль (67.8--100.4 кДж/моль). Однако эти значения были получены для катализатора Rh/Al$_{2}$O$_{3}$, в случае же применения оксида Ce0.75Zr0.25O2 в качестве носителя следует ожидать более низкие значения, вследствие наличия в этом оксиде кислородных вакансий, которые сильнее активируют диоксид углерода в процессе метанирования CO$_{2}$. Для процессов паровой конверсии метана при использовании родиевых катализаторов наблюдаемая энергия активации выше, чем в паровом риформинге более тяжелых углеводородов. Это связано с более прочной связью C-H в молекуле метана и ее низкой поляризации.

В работах \cite{Schadel2009,Zeppieri2010,Jakobsen2010} было показано, что для родиевых катализаторов энергия активации паровой конверсии метана сильно зависит от условий проведения реакции и состава используемого катализатора и варьируется от 40 до 90 кДж/моль.

Таким образом, для родиевых катализаторов наблюдаемые энергии активации в процессах конверсии пропана и этана обычно составляют 20--70 кДж/моль, конверсия метана характеризуется более высокими значениями энергии активации (40--90 кДж/моль). Для процесса метанирования диоксида углерода --- 67.8--100.4 кДж/моль. Энергии активации, полученные в нашем исследовании, соответствуют литературным данным, что указывает на корректность используемого приближения.

Расчетные парциальные давления были пересчитаны с учетом $k_0$ и $E_a$, полученных методом наименьших квадратов. Было рассчитано среднее (по всем температурам) абсолютное отклонение расчетных парциальных давлений от экспериментальных для каждого вещества по формуле

\begin{equation}\label{eq15}
\varepsilon_{i}=\frac{\sum_{j=1}^{T_{count}}\left|\overline{P_{ij}}- \widetilde{P_{ij}} \right|}{T_{count}}, i \in [1, n],
\end{equation}
где $\overline{P_{ij}}$, $\widetilde{P_{ij}}$ --- экспериментальное и расчетное значения парциальных давлений вещества $i$ для температуры $j$, атм; $T_{count}$ --- количество температур; $n$ --- количество измеряемых веществ. В таблице ~\ref{tab:abs_error} представлены средние значения абсолютных отклонений расчетных парциальных давлений от экспериментальных, а также максимальные значения экспериментальных парциальных давлений соответствующих веществ.
\begin{table}[h]
\noindent
\caption{Средние абсолютные отклонения расчетных парциальных давлений от экспериментальных и максимальные экспериментальные парциальные давления}
\label{tab:abs_error}
\centering
\begin{tabular}{|c|c|c|c|c|c|c|}
\hline
\textbf{Вещество} & \textbf{$\varepsilon_{i}$, атм} & \textbf{$\overline{P_{max}}$, атм} \\
\hline
\ce{C3H8}  & 0.009   & 0.060 \\
\ce{H2O}      & 0.013   & 0.370 \\
\ce{CO2}      & 0.010   & 0.049 \\
\ce{H2}       & 0.015   & 0.079 \\
\ce{CH4}      & 0.062   & 0.626 \\
\ce{C2H6}  & 0.006   & 0.030 \\
\hline
\end{tabular}
\end{table}

Статистические характеристики можно считать хорошими для \ce{H2O} и \ce{CH4}, но недостаточными для веществ \ce{C3H8}, \ce{CO2}, \ce{H2}, \ce{C2H6}.
На рис.~\ref{fig:result} представлены графики зависимости экспериментальных и расчетных парциальных давлений от температуры для \ce{C3H8}, \ce{H2O}, \ce{CO2}, \ce{H2}, \ce{CH4}, \ce{C2H6}. 
Экспериментальные данные обозначены синим цветом, а расчетные --- оранжевым. Полученные результаты подтверждают, что модель соответствует эксперименту и может использоваться для прогнозирования состава продуктов паровой конверсии углеводородов.

Качественное соответствие модели экспериментальным данным является приемлемым. Однако количественная аппроксимация в настоящий момент не достаточна для веществ \ce{C3H8}, \ce{CO2}, \ce{H2}, \ce{C2H6}. В продолжение работы предполагается включить в модель дополнительные параметры для более точного описания зависимостей.


%\clearpage
\begin{figure}[ht]
    \centering
    \begin{subfigure}[b]{0.45\linewidth}
        \centering
        \includegraphics[width=\linewidth]{C3H8.jpg}
        \caption{\ce{C3H8}}
    \end{subfigure}
    \hfill
    \begin{subfigure}[b]{0.45\linewidth}
        \centering
        \includegraphics[width=\linewidth]{H2O.jpg}
        \caption{\ce{H2O}}
    \end{subfigure}

    \begin{subfigure}[b]{0.45\linewidth}
        \centering
        \includegraphics[width=\linewidth]{CO2.jpg}
        \caption{\ce{CO2}}
    \end{subfigure}
    \hfill
    \begin{subfigure}[b]{0.45\linewidth}
        \centering
        \includegraphics[width=\linewidth]{H2.jpg}
        \caption{\ce{H2}}
    \end{subfigure}

    \begin{subfigure}[b]{0.45\linewidth}
        \centering
        \includegraphics[width=\linewidth]{CH4.jpg}
        \caption{\ce{CH4}}
    \end{subfigure}
    \hfill
    \begin{subfigure}[b]{0.45\linewidth}
        \centering
        \includegraphics[width=\linewidth]{C2H6.jpg}
        \caption{\ce{C2H6}}
    \end{subfigure}
    \caption{Сопоставление экспериментальных и расчетных парциальных давлений}\label{fig:result}
\end{figure}

\section*{Заключение}
\label{sec-conclusion}

Результатом проведенного исследования является решение задачи выбора кинетических параметров для процесса низкотемпературной паровой конверсии пропана на 1~вес.\,\%~Rh/Ce$_{0.4}$Zr$_{0.5}$Y$_{0.05}$La$_{0.05}$O$_{2}$ катализаторе. Найденные кинетические параметры будут использоваться для получения газовых смесей с необходимыми характеристиками.
Также в дальнейших исследованиях будет проведено усложнение полученной модели для учета реакций паровой конверсии более тяжелых углеводородов (бутана, пентана), входящих в состав ПНГ. Это позволит увеличить диапазон условий применимости модели и заранее оценивать возможность применения процесса НТПК для эффективной утилизации ПНГ на различных месторождениях.

В планы дальнейших работ входит использование других схем распараллеливания, способных повысить эффективность использования параллельных вычислительных систем для решения задач глобальной оптимизации. Сложности с достижением значительных показателей ускорения объясняются следующими двумя обстоятельствами. 

Во-первых, при реализации параллельного алгоритма была применена синхронная схема распараллеливания. Управляющий поток-мастер (в котором работал метод оптимизации) при выполнении  шага 5 алгоритма ожидал окончания проведения всех испытаний, проводимых рабочими потоками. В решаемой задаче время проведения одного испытания зависит от времени решения системы ОДУ и может варьироваться в зависимости от конкретных значений параметров. Возможный путь к повышению эффективности распараллеливания --- реализация асинхронной схемы проведения испытаний на шаге 5 алгоритма. В этом случае результаты испытания, полученные одним из вычислителей, могут быть отправлены для обработки мастеру немедленно, без ожидания окончания других испытаний. После чего вычислитель получает от мастера новую точку, в которой начнет проведение очередного испытания. Асинхронная схема распараллеливания, теоретически, способна обеспечить полную загрузку задействованных вычислительных ресурсов даже в том случае, когда время выполнения испытаний различно в разных точках области поиска.
 
Во-вторых, при масштабном распараллеливании <<узким местом>> становится локальный поиск, который выполняется для уточнения решения после завершения фазы глобального поиска. Локальный поиск требует незначительного (по сравнению с глобальным) числа поисковых испытаний, но эти испытания выполняются последовательно. И при параллельном запуске это незначительное число последовательных испытаний становится сопоставимым с числом параллельных итераций глобального поиска. В дальнейшем планируется использовать параллельную реализацию метода локального поиска на завершающем этапе работы, что способно повысить эффективность распараллеливания в целом.

Наконец, планируется развитие математических основ используемых алгоритмов параллельной глобальной оптимизации. 
В задаче поиска параметров модели, обеспечивающих минимум невязки между расчетными и экспериментальными данными, теоретическое значение целевой функции в точке глобального минимума является известным --- оно равно нулю.
При этом точка глобального минимума (соответствующая лучшим значениям параметров модели) является неизвестной и подлежит определению.
На практике нулевое значение невязки может и не достигаться, если исследуемая модель плохо соответствует эксперименту, но в таком случае целевая функция будет положительной.  
Информацию о нижней границе значений целевой функции можно использовать для разработки нового алгоритма оптимизации, который позволит как гарантировать сходимость к глобальному минимуму, так и ускорить процесс его отыскания.

\vspace{1em}
{\it Работа выполнена при поддержке Министерства образования и науки Российской Федерации, проекты № FSWR-2026-0007 (разработка параллельного алгоритма и проведение вычислительных экспериментов, К.А. Баркалов, И.Г. Лебедев) и № FMRS-2025-0031 (разработка кинетической модели, И.М. Губайдуллин).}


\printbibliography[title={СПИСОК ЛИТЕРАТУРЫ}]

% Сведения об авторах (Ф.И.О.(полностью), степень, звание, кафедра (отдел), вуз, e-mail)
Баркалов Константин Александрович, д.т.н., доцент, заведующий кафедрой математического обеспечения и суперкомпьютерных технологий, Нижегородский государственный университет им. Н.И. Лобачевского (Н. Новгород, Российская Федерация)

Варсеева Екатерина Николаевна, ассистент кафедры математического обеспечения и суперкомпьютерных технологий, Нижегородский государственный университет им. Н.И. Лобачевского (Н. Новгород, Российская Федерация)

Лебедев Илья Генадьевич, к.т.н., доцент кафедры математического обеспечения и суперкомпьютерных технологий, Нижегородский государственный университет им. Н.И. Лобачевского (Н. Новгород, Российская Федерация)

Хашпер Анна Леонидовна, главный специалист отдела бизнес-систем, ООО <<СИБИНТЕК-СОФТ>> (Уфа, Российская Федерация)

Урлуков Артем Сергеевич, аспирант, младший научный сотрудник отдела гетерогенного катализа, Институт катализа им. Г.К. Борескова СО РАН
(Новосибирск, Российская Федерация)

Губайдуллин Ирек Марсович, профессор, д.ф.-м.н., заведующий лабораторией математической химии ИНК УФИЦ РАН, профессор кафедры технологии нефти и газа УГНТУ
(Уфа, Российская Федерация)

Хисаметдинов Искандер Фанзилевич, бакалавр Евразийской политехнической школы УГНТУ (Уфа, Российская Федерация)


\newpage
{\phantom{eee} \hrule \vskip 4 pt}
\classify{} % Код MSC не нужен
%\classify{MSC 00000} % Коды MSC см. http://www.ams.org/msc/

\title{PARALLEL ALGORITHM OF GLOBAL OPTIMIZATION AND ITS USE FOR SOLVING INVERSE PROBLEMS OF CHEMICAL KINETICS}

\authorsen{K.A.~Barkalov$^{1,2}$, E.N.~Varseeva$^{1}$, I.G.~Lebedev$^{1}$, A.L.~Khashper$^{3}$, A.S.~Urlukov$^{4,5}$, I.M.~Gubaydullin$^{6}$, I.F.~Khisametdinov$^{7}$}

\organizations{
$^{1}$Lobachevsky State University of Nizhny Novgorod \\ (pr. Gagarina 23, N. Novgorod, 603022 Russia)\\
$^{2}$Scientific and Educational Mathematical Center <<Mathematics of Future Technologies>> \\ (pr. Gagarina 23 bldg.2, N. Novgorod, 603022 Russia) \\
$^{3}$LLC <<SIBINTEK-SOFT>> \\ (Kommunalnyy Proezd 30, Khimki, 141401 Russia)\\
$^{4}$Catalysis Institute, Siberian Branch of the Russian Academy of Sciences \\ (pr. Academician Lavrentiev 5, Novosibirsk, 630090 Russia)\\
$^{5}$Novosibirsk State University \\ (Pyrogova Street 1, Novosibirsk, 630090 Russia) \\
$^{6}$Institute of Petrochemistry and Catalysis %--- Subdivision of Ufa Federal Research Centre of the RAS 
\\ (pr. Oktyabrsky 141, Ufa, 450075 Russia)\\
$^{7}$Ufa State Petroleum Technological University \\ (Kosmonavtov Street 1, Ufa, 450064 Russia)\\
}
\emails{\emailsinsert}
\redactiondateen{\rdate}

\begin{abstract}%
The article discusses the use of parallel computing for selecting parameters of a mathematical model of low-temperature steam conversion process of hydrocarbons contained in associated petroleum gas. For this chemical process, it is necessary to develop its kinetic model, i.e., determine the corresponding kinetic reaction parameters. To achieve this, an inverse problem is solved where values of kinetic parameters are sought based on experimental data. Mathematically, the inverse problem of chemical kinetics corresponds to a global optimization problem. A parallel information-statistical algorithm for global search combined with local refinement of the best solution was used to solve this problem. 
The algorithm is deterministic and relies on the assumption of Lipschitz continuity of the objective function, which is typical for many other approaches to constructing global optimization methods. In this case, solving multidimensional problems reduces to solving equivalent one-dimensional problems. The corresponding reduction is based on using a Peano curve that maps the unit interval of the real axis onto a hypercube. Parallelization of the algorithm is organized according to the <<master-workers>> scheme with shared memory usage.
The found optimal model parameters allowed adequately simulating the process of low-temperature steam conversion of light hydrocarbons using a catalyst.

\keywordsen{global optimization, multi-extremal functions, parallel computing, chemical kinetics, inverse problems}
\end{abstract}

\forcitationen{Barkalov~K.A., Varseeva~E.N., Lebedev~I.G. Parallel algorithm of global optimization and its use for solving inverse problems of chemical kinetics}

\openaccess



% =============== СПИСОК "REFERENCES" (английская часть) ===============

\begin{refsection}[references_en.bib]
\nocite{*}
\printbibliography[title={REFERENCES}]
\end{refsection}
\end{document}